% ============================================================
% APPENDIX: GLOSSARY OF TERMS
% ToggleTails
% ============================================================

\chapter*{Appendix A: Glossary of Terms}
\addcontentsline{toc}{chapter}{Appendix A: Glossary of Terms}
\label{appendix:glossary}

The following glossary defines technical, domain-specific, and
application-specific terminology used throughout this report. Terms are
listed alphabetically within thematic groupings.

% ─────────────────────────────────────────────────────────────────────────────
\section*{A.1\quad Mobile Development}
% ─────────────────────────────────────────────────────────────────────────────

\begin{description}[leftmargin=2.5cm, style=nextline]

\item[\textbf{Android Keystore}]
A platform-native secure storage mechanism on Android devices that provides
hardware-backed encryption for sensitive credentials. Used in ToggleTails
via \texttt{expo-secure-store} to protect the parent PIN hash.

\item[\textbf{AsyncStorage}]
A key-value local storage system provided by the
\texttt{@react-native-async-storage/async-storage} library for React Native.
Used in ToggleTails to persist child profiles, story content, approval states,
reading progress, and narration settings across sessions. Unlike
\texttt{expo-secure-store}, AsyncStorage does not provide hardware-level
encryption and is not suitable for storing sensitive credentials.

\item[\textbf{Cross-Platform Development}]
A software development approach in which a single codebase compiles and runs
on multiple operating systems. ToggleTails uses React Native with Expo to
target both iOS and Android from one TypeScript codebase.

\item[\textbf{EAS (Expo Application Services)}]
A suite of cloud-based build and deployment tools provided by Expo for
compiling native iOS and Android application binaries from a React Native
codebase. EAS is used in ToggleTails to produce native builds required for
features such as speech recognition, and to generate production-ready
artifacts for app store distribution.

\item[\textbf{Expo}]
An open-source framework and platform built on top of React Native that
simplifies cross-platform mobile development by providing a managed build
system, a standard library of device APIs, and the Expo Go development client.
ToggleTails is built using Expo SDK 54.

\item[\textbf{Expo Go}]
A development client application provided by Expo that allows React Native
applications to be previewed on physical devices without a native build. Some
features—specifically speech recognition—are unavailable in Expo Go and
require an EAS native build.

\item[\textbf{Expo Router}]
A file-based navigation library for React Native built on top of React
Navigation. In ToggleTails, screen files placed in the \texttt{app/}
directory are automatically mapped to application routes, supporting both
stack and tab navigation patterns.

\item[\textbf{expo-av}]
An Expo library providing audio recording and playback capabilities. Used in
ToggleTails for recording parent narration audio and for playing back both
human recordings and ElevenLabs-synthesised audio.

\item[\textbf{expo-secure-store}]
An Expo library that stores key-value data using platform-native secure
enclave mechanisms: the iOS Keychain on Apple devices and the Android
Keystore on Android devices. Used in ToggleTails to store the parent PIN hash.

\item[\textbf{expo-speech}]
An Expo library providing on-device text-to-speech synthesis. Used in
ToggleTails as the offline fallback narration engine when ElevenLabs is
unavailable.

\item[\textbf{expo-speech-recognition}]
An Expo library enabling continuous speech recognition using native device
microphone input. Used in the Help Me Read feature to capture and evaluate
the child's spoken reading. Requires a native EAS build.

\item[\textbf{iOS Keychain}]
A platform-native secure storage mechanism on iOS and macOS devices that
provides encrypted storage for sensitive application data. Used in ToggleTails
via \texttt{expo-secure-store} to protect the parent PIN hash.

\item[\textbf{Metro Bundler}]
The JavaScript bundler used by React Native and Expo to package application
source files and assets into a format suitable for execution on mobile
devices. Metro handles dynamic imports of the bundled JSON story files in
ToggleTails.

\item[\textbf{React Native}]
An open-source mobile application framework maintained by Meta that enables
the development of native iOS and Android applications using JavaScript and
React. ToggleTails is built on React Native v0.81.5.

\item[\textbf{TypeScript}]
A strongly typed superset of JavaScript that adds static type checking at
compile time. Used throughout the ToggleTails mobile codebase to improve
reliability, tooling, and code documentation.

\end{description}

% ─────────────────────────────────────────────────────────────────────────────
\section*{A.2\quad Backend and Data}
% ─────────────────────────────────────────────────────────────────────────────

\begin{description}[leftmargin=2.5cm, style=nextline]

\item[\textbf{API (Application Programming Interface)}]
A defined contract through which software components communicate. In
ToggleTails, the backend exposes a RESTful HTTP API that the mobile client
calls for story generation, ElevenLabs narration, child profile management,
approvals, and synchronisation.

\item[\textbf{bcrypt / bcryptjs}]
A password-hashing algorithm that applies a configurable cost factor and a
random salt to produce a one-way hash. Used in ToggleTails to hash parent
passwords and PINs before storage, providing resistance to brute-force and
rainbow-table attacks.

\item[\textbf{Bearer Token}]
An HTTP authentication scheme in which a signed token is transmitted in the
\texttt{Authorization} header of a request. ToggleTails uses Bearer-scheme
JWTs to authorise parent access to protected backend endpoints.

\item[\textbf{CORS (Cross-Origin Resource Sharing)}]
An HTTP mechanism that controls which origins are permitted to make requests
to a server. The ToggleTails backend enables CORS globally for development;
origin restriction would be applied for production deployment.

\item[\textbf{CRUD}]
An acronym for the four basic persistent storage operations: Create, Read,
Update, and Delete. The ToggleTails backend implements CRUD operations for
parent accounts, child profiles, story records, and approval records.

\item[\textbf{dotenv}]
A Node.js library that loads environment variables from a \texttt{.env} file
into \texttt{process.env} at runtime. Used in the ToggleTails backend to
manage API keys, database credentials, and JWT secrets without hardcoding
them in source files.

\item[\textbf{Express.js}]
A minimal and flexible Node.js web application framework used to build the
ToggleTails backend REST API. Express handles routing, middleware composition,
and HTTP request/response processing.

\item[\textbf{HTTP Status Codes}]
Standardised numeric codes returned by HTTP servers to indicate the outcome
of a request. Codes in the 2xx range indicate success; 4xx indicate client
errors (such as 401 Unauthorized or 404 Not Found); 5xx indicate server
errors.

\item[\textbf{JSON (JavaScript Object Notation)}]
A lightweight, human-readable data interchange format used throughout
ToggleTails for API request and response payloads, story content files,
and the story manifest.

\item[\textbf{JWT (JSON Web Token)}]
A compact, self-contained token format used for stateless authentication. A
JWT encodes a payload (such as a parent identifier) and an expiry time,
signed with a secret key. The ToggleTails backend issues JWTs upon successful
authentication and validates them on protected routes.

\item[\textbf{Middleware}]
A function in a server-side framework that executes in the request-handling
pipeline before the final route handler. In ToggleTails, middleware is used
for JWT authentication enforcement (\texttt{requireAuth}), global error
handling (\texttt{errorHandler.js}), and CORS configuration.

\item[\textbf{MySQL}]
An open-source relational database management system. The ToggleTails backend
uses MySQL as its persistence layer, accessed via the Prisma ORM.

\item[\textbf{Node.js}]
A JavaScript runtime built on the V8 engine that enables server-side
JavaScript execution. The ToggleTails backend API is implemented in Node.js.

\item[\textbf{ORM (Object-Relational Mapping)}]
A programming technique that maps database tables to objects in application
code, abstracting raw SQL. ToggleTails uses Prisma as its ORM, providing a
type-safe query API over MySQL.

\item[\textbf{Prisma}]
A modern Node.js ORM that generates a type-safe client from a declarative
schema file (\texttt{schema.prisma}). Prisma handles database migrations,
query construction, and relation management in the ToggleTails backend.

\item[\textbf{REST (Representational State Transfer)}]
An architectural style for distributed systems in which resources are
identified by URIs and accessed using standard HTTP methods (GET, POST, PATCH,
DELETE). The ToggleTails backend API follows REST conventions.

\item[\textbf{Salt (Cryptographic)}]
A random value added to input data before hashing to ensure that identical
inputs produce different hashes. bcryptjs generates a unique salt for each
credential hash, preventing precomputed rainbow-table attacks.

\item[\textbf{Upsert}]
A database operation that inserts a new record if one does not already exist,
or updates the existing record if it does. Used in ToggleTails for story
approval records, ensuring one record per child--story pair without
duplicates.

\end{description}

% ─────────────────────────────────────────────────────────────────────────────
\section*{A.3\quad Artificial Intelligence and Narration}
% ─────────────────────────────────────────────────────────────────────────────

\begin{description}[leftmargin=2.5cm, style=nextline]

\item[\textbf{ElevenLabs}]
A commercial neural text-to-speech and voice cloning platform. ToggleTails
uses the ElevenLabs API to generate high-quality AI narration audio and,
optionally, to clone a parent's voice from a recorded sample for use in AI
narration sessions.

\item[\textbf{Fuzzy Matching}]
A string comparison technique that identifies approximate matches between
two strings, tolerating minor differences such as pronunciation variation or
transcription noise. Used in the Help Me Read feature to compare the child's
spoken words against expected page text.

\item[\textbf{GPT-4o-mini}]
A compact variant of OpenAI's GPT-4o large language model. Used in
ToggleTails for AI story generation, producing age-appropriate, personalised
narrative text from structured prompts encoding child profile attributes.

\item[\textbf{Large Language Model (LLM)}]
A machine learning model trained on large corpora of text data, capable of
generating coherent and contextually relevant natural language output. The
ToggleTails story generation pipeline uses OpenAI's GPT-4o-mini as its LLM.

\item[\textbf{Neural Text-to-Speech (Neural TTS)}]
A class of speech synthesis systems that use deep neural networks to produce
natural-sounding audio from written text. Neural TTS systems such as
ElevenLabs produce significantly more natural prosody and voice quality than
older rule-based or concatenative TTS systems.

\item[\textbf{OpenAI}]
An AI research and deployment company. ToggleTails uses the OpenAI API to
access the GPT-4o-mini model for AI story generation.

\item[\textbf{Prompt Engineering}]
The practice of designing and structuring input text (prompts) to guide the
output of a large language model. In ToggleTails, the story generation backend
constructs prompts that encode child age, reading level, interests, and
optional custom titles to produce well-formed, age-appropriate stories.

\item[\textbf{Speech Recognition}]
The process of converting spoken audio into text. ToggleTails uses
\texttt{expo-speech-recognition} in the Help Me Read feature to transcribe the
child's spoken reading in real time for comparison against expected page text.

\item[\textbf{Text-to-Speech (TTS)}]
The process of synthesising spoken audio from written text. ToggleTails
implements TTS through two mechanisms: on-device synthesis via
\texttt{expo-speech}, and neural synthesis via the ElevenLabs API.

\item[\textbf{Voice Cloning}]
A process in which a synthetic voice model is trained from one or more audio
samples of a real person's voice, enabling AI synthesis to approximate that
person's vocal characteristics. ToggleTails optionally submits parent
recordings to the ElevenLabs Voice Cloning API to create a cloned voice for
AI narration.

\end{description}

% ─────────────────────────────────────────────────────────────────────────────
\section*{A.4\quad System Architecture and Design}
% ─────────────────────────────────────────────────────────────────────────────

\begin{description}[leftmargin=2.5cm, style=nextline]

\item[\textbf{Client--Server Architecture}]
A distributed system model in which client applications (such as the
ToggleTails mobile app) send requests to a centralised server (the Node.js
backend) that processes them and returns responses.

\item[\textbf{Graceful Degradation}]
A design principle in which a system maintains reduced but functional
operation when a component fails or is unavailable. ToggleTails implements
graceful degradation throughout: the backend is optional for core reading;
ElevenLabs falls back to \texttt{expo-speech}; the database failure allows
partial server operation.

\item[\textbf{Layered Architecture}]
An architectural pattern in which a system is divided into horizontal layers,
each with a distinct responsibility, communicating only with adjacent layers.
ToggleTails applies a four-layer model within the mobile client: Presentation,
Domain, Data, and Platform Services.

\item[\textbf{Local-First Architecture}]
A design philosophy in which application data is stored and processed
locally on the user's device as the primary source of truth, with network
synchronisation as an optional enhancement rather than a requirement.
ToggleTails adopts a local-first approach to ensure offline availability and
protect user privacy.

\item[\textbf{Modular Design}]
A software engineering principle in which a system is decomposed into
self-contained, independently maintainable components (modules), each with a
clearly defined responsibility and interface. ToggleTails separates
narration, storage, authentication, approval, and seeding into distinct
service modules.

\item[\textbf{Offline-First}]
A software design strategy in which an application is designed to function
without an internet connection as the baseline, with online features layered
on top. See also: Local-First Architecture.

\item[\textbf{Role-Based Access Control (RBAC)}]
A security model in which access to system features is determined by the role
assigned to the current user. ToggleTails implements two roles: Parent
(authenticated, full configuration access) and Child (unauthenticated,
read-only story access).

\item[\textbf{Separation of Concerns}]
A software design principle in which different responsibilities are handled
by distinct components, minimising coupling and improving maintainability.
Reflected in ToggleTails through the separation of presentation, domain,
data, and platform layers, and through the separation of Parent View and
Child View.

\item[\textbf{Waterfall Methodology}]
A sequential software development model in which each phase (requirements,
design, implementation, testing, evaluation) must be completed before the
next begins. ToggleTails is developed using the Waterfall model across four
defined phases.

\end{description}

% ─────────────────────────────────────────────────────────────────────────────
\section*{A.5\quad Application-Specific Terms}
% ─────────────────────────────────────────────────────────────────────────────

\begin{description}[leftmargin=2.5cm, style=nextline]

\item[\textbf{AI Narration Mode}]
A narration mode in ToggleTails in which story audio is generated by the
ElevenLabs neural TTS API (or, as a fallback, by \texttt{expo-speech}).
Contrast with Human Narration Mode.

\item[\textbf{Child View}]
The simplified, unauthenticated interface in ToggleTails through which the
child browses stories, selects narration, and tracks reading progress.
Configuration features are intentionally absent from this view.

\item[\textbf{Help Me Read}]
An interactive reading mode in ToggleTails in which the application listens
to the child read aloud and provides real-time feedback on pronunciation
accuracy using speech recognition and fuzzy word matching.

\item[\textbf{Human Narration Mode}]
A narration mode in ToggleTails in which a story is narrated using audio
recorded by the parent. The recording is stored locally and played back via
\texttt{expo-av} during the child's reading session.

\item[\textbf{Parent Gate}]
The PIN-based authentication mechanism in ToggleTails that restricts access
to the Parent View. The parent creates a PIN during onboarding; it is hashed
with bcrypt and stored in \texttt{expo-secure-store}.

\item[\textbf{Parent View}]
The authenticated configuration interface in ToggleTails accessible only to
guardians after PIN verification. Provides access to story approval, narration
settings, child profile management, and AI story creation.

\item[\textbf{PreloadedSeeder}]
A module in ToggleTails (\texttt{src/data/seeder/PreloadedSeeder.ts}) that
seeds the 52 bundled story JSON files into AsyncStorage on first application
launch, marking all pre-loaded stories as approved by default.

\item[\textbf{Reading Level}]
A classification used in ToggleTails to match story complexity to the child's
literacy stage. Three levels are defined: Beginner, Intermediate, and
Advanced. The reading level is set by the parent during onboarding and stored
in the child profile.

\item[\textbf{Story Approval Workflow}]
The process in ToggleTails by which a parent reviews and approves
AI-generated stories before they become visible to the child. Pre-loaded
stories are auto-approved on first seed; AI-generated stories default to
\texttt{approved: false} until the parent acts.

\item[\textbf{Story Library}]
The curated collection of 52 pre-loaded public-domain stories bundled with
the ToggleTails application, organised across six genres (Animals, Adventure,
Bedtime, Science, Values, Fantasy) and three reading levels. Stories are
defined in JSON and loaded via \texttt{storyLibraryService.ts}.

\item[\textbf{storyLibraryService}]
A TypeScript module (\texttt{src/data/library/storyLibraryService.ts}) that
provides the primary API for accessing the pre-loaded story library. Exposes
functions for genre-filtered retrieval, reading-level filtering, paginated
listing, and lazy full-content loading from the bundled JSON manifest.

\item[\textbf{ToggleTails}]
The mobile storytelling application developed as the engineering artefact of
this capstone project. ToggleTails provides a child-centred reading experience
with a dual-mode narration system (human and AI narration) governed by a
parent-controlled configuration interface.

\end{description}

% ─────────────────────────────────────────────────────────────────────────────
\section*{A.6\quad Accessibility and Standards}
% ─────────────────────────────────────────────────────────────────────────────

\begin{description}[leftmargin=2.5cm, style=nextline]

\item[\textbf{HTTPS (Hypertext Transfer Protocol Secure)}]
An extension of HTTP that uses TLS encryption to secure data in transit
between client and server. The ToggleTails backend operates over HTTP during
development; HTTPS would be enforced via a reverse proxy or managed cloud
service prior to production deployment.

\item[\textbf{Informed Consent}]
A process by which research participants (or, in the case of minors, their
legal guardians) are provided with sufficient information about a study to
make a voluntary and informed decision to participate. All observational
sessions in the ToggleTails user study were conducted under informed parental
consent.

\item[\textbf{WCAG (Web Content Accessibility Guidelines)}]
An internationally recognised set of guidelines published by the W3C for
making digital content accessible to people with disabilities. ToggleTails
targets WCAG compliance for colour contrast ratios and touch target sizes to
support users with varying visual and motor abilities.

\end{description}
